\chapter{Fourier Transform and Laplace Transform}
\renewcommand{\currentchapterpath}{chapters/ft/}

\begin{flushright}
\textit{“往来不穷谓之通。” ——《周易·系辞上》}
\footnote{From the ``Appended Statements'' of the \textit{Book of Changes}: ``Coming and going without exhaustion is called penetration.'' An integral transform sends a function into another domain and the inverse returns it. Taken together, the two maps make a problem transparent: differentiation becomes multiplication, convolution becomes a product, and an initial-value problem becomes algebra.}
\end{flushright}

The Fourier transform extends a trigonometric expansion from periodic functions to functions on the real line. The Laplace transform is a closely related tool for causal functions and initial-value problems. This chapter passes from Fourier series to the Fourier integral, develops the delta function as the natural language of point sources and inversion, and then introduces the Laplace transform.

\chapterinput{fourier_transform}
\chapterinput{delta_function}
\chapterinput{laplace_transform}
\chapterinput{applications}
\chapterinput{exercises}
